% ============================================================================
%  RESULTS -- MAIN TEXT, SECTION 3: ESTIMATION
%  project `panel-npiv-discrete`
%
%  WHAT THIS FILE HOLDS.  The 4 statements that belong under Section 3:
%  the estimator in all four of its pieces: the sieve Landweber--Fridman
%  varphi-step, the debiased moment function with its two matrices, the
%  cross-fitting scheme, and the preliminary and one-step theta-estimator.
%  Nothing else.  In order: def:lfstep, def:moment, def:crossfit, def:onestep.
%
%  WHY THREE FILES AND NOT ONE.  results_main.tex printed all 23 core statements
%  in one block, which the article had to \input somewhere; it went under
%  Section 4, so Section 2's prose forward-referenced statements forty pages
%  later.  The three files are \input under the three sections that discuss
%  them.  The split is a PARTITION: results_s2 + results_s3 + results_s4 +
%  results_supp = the undivided results.tex, every \label declared exactly once
%  across the four, no statement rewritten, retitled or renumbered.  Relative
%  order inside each file is the order of the parent file.
%
%%  DO NOT \input THIS FILE ALONGSIDE results.tex OR assumptions.tex.
%%  results.tex and assumptions.tex SURVIVE UNDIVIDED and remain the source of
%%  record: scripts/build-theory/build.sh compiles them, scripts/render-sync/
%%  check_rendering_sync.py reads them, and ~100 files under scripts/symbolic/
%%  specs/ cite them BY LINE NUMBER.  Neither was edited.  Every \label in this
%%  file is therefore declared TWICE on disk -- here and in its parent -- and the
%%  two must never reach the same document.
%
%  SUPPLEMENT POINTERS.  A statement kept here may cite a statement that lives
%  in the supplement.  Those calls are written \suppref{Kind}{name} ->
%  "Kind S:name" and \suppeq{name} -> "(S:name)", both defined in
%  paper_header.tex.  They are not \cref calls and cannot render "??".
%
%  WHAT def:lfstep AND def:crossfit NO LONGER CARRY.  Those two ran 286 and 321
%  live lines, about 17 pages between them, on the same triage the assumption
%  block took (1090 -> 404 live lines).  What stays is the CONSTRUCTION, and it
%  stays verbatim: the objects, the recursion, the stopping rule, the fold
%  split, and every clause tag any other file cites.  What moved to
%  results_s3_supp.tex is the measurability verifications, the "why this
%  container and not that one" paragraphs, the re-declarations of conventions
%  stated elsewhere, and the discussion of what each clause rules out.  Each
%  moved block left a sentence here saying the condition holds and pointing at
%  "the supplement's discussion of the estimator", which is that file's
%  \section*{} heading.
%
%  NO LABELLED DISPLAY MOVED, so no \eqref anywhere was orphaned and no equation
%  number changed: eq:normconst, eq:cap, eq:lfsched, eq:prodmeas,
%  eq:supnormconv and eq:foldrates all stayed, and results_s3_supp.tex declares
%  no \label and contains no numbered environment.
%
%  THREE THINGS WERE KEPT DESPITE LOOKING LIKE COMMENTARY, because article-
%  visible text cites them:  the sentence fixing cal I_star (P-a.1) in
%  results_s4 reads "the fold cal I_star of def:lfstep";  (c'.2)'s [0,+infty]
%  typing, which (R-a) in results_s4 named;  and (d)'s "strengthening of A5(iv)
%  in form", which (R-e.1) cited as "the reason def:crossfit(d) gives".
%
%  AMENDED 2026-08-02.  The second and third citers moved out of the STATEMENTS
%  of thm:rootn and into rem:rootngroups(3) and rem:rootngroups(7), in the
%  statement-shortening pass that enforced .claude/references/theorem-
%  environment-style.md.  BOTH ARE STILL CITED, from the supplement rather than
%  from the article body, so the three items above are still load-bearing and
%  none of them may be moved out of this file.
% ============================================================================


\begin{definition}[The estimator: iterated linearisation with sieve Landweber--Fridman]\label{def:lfstep}
Let $\hat\varphi^{(0)}$ be a preliminary estimator which is a \emph{measurable random element of
$\mathcal V$ with respect to the sample it is computed on} \textbf{and which lies in $\mathcal S_n$}, the
sieve fixed at the foot of this paragraph; and let $\hat\varepsilon_0\to0$ be a \emph{known}
deterministic sequence with
$P\big(\|\hat\varphi^{(0)}-\vdagger\|_{L^{2}(\pi)}\le\hat\varepsilon_0\big)\to1$ and
$P\big(\|\hat\varphi^{(0)}-\vdagger\|_{H^{s}}\le\mathsf R\big)\to1$, the second being a statement about
a finite quantity only because $\hat\varphi^{(0)}\in\mathcal S_n$ and \eqref{eq:sieveinhs} places
$\mathcal S_n$ in $H^{s}$, while $\vdagger\in H^{s}$ by H1's source set.
Let
$\mathcal S_n:=\operatorname{span}(b_1,\dots,b_{L_x})\cap\mathcal V$ with $b_1,\dots,b_{L_x}$
the B-spline basis of A8. \emph{Normalise the scale.} Let
\begin{equation}\label{eq:normconst}
\omega_n\;:=\;2\,(2c_C)^{1/2}\,\mathsf c_2^{2}\,c_\Pi\,\|1/\pi\|_\infty^{1/2}\,h_u^{-1}h_x^{-1/2}
\end{equation}
% c_2 GLOSSED HERE 2026-08-05: it sits inside eq:normconst, a numbered display the
% article prints, and was defined only in the supplement.
where $\mathsf c_2:=\|\bar\Theta\|_\infty<\infty$ is the sup-norm of C3(ii)'s kernel envelope and
$c_C:=\|\bar\Theta\|_{L^{1}}/(2\mathsf c_2)$ is the \emph{envelope width constant}, so that
$2c_C\mathsf c_2=\|\bar\Theta\|_{L^{1}}$; both are fixed by C3(ii), and $c_\Pi$ is the range-measure
constant of B1(iv.c).
% c_2 GLOSSED HERE 2026-08-05; RE-POINTED AT THE ENVELOPE 2026-08-15.  Through
% dispatch 25 c_C was the kernel's SUPPORT HALF-WIDTH and c_2 the sup of
% Theta_{s_0} over |varsigma| <= c_C.  The shipped kernel k_order(., 8) is
% Gaussian-based and has NO compact support, so no finite c_C made that reading
% true.  C3(ii) now posits an ENVELOPE and DERIVES c_C from it; eq:normconst is
% unchanged glyph for glyph, and 2 c_C c_2 is the same number it always was --
% the L^1 mass of the envelope.  Measured at p_K = 8 AT THE MINIMAL ENVELOPE THAT
% lem:bdryfam(d) FIXES: c_2 = 34.0, c_C = 0.155, 2 c_C c_2 = 10.56 -- against the
% retired half-width reading's ~7, so the display keeps its numerical meaning and
% not merely its spelling.  A d27 draft of this comment said c_C = 1.82, which is
% the ANALYTIC MAJORANT's value and pairs with ITS c_2 = 4.50e8, not with anything
% eq:normconst is instantiated at; eq:envconst binds the two through
% int bar-Theta = 2 c_C c_2, so a c_C off one envelope and a c_2 off another
% describe no envelope at all.
This is
the bound of \suppref{Proposition}{hatAprop}(i), which is \emph{deterministic} and holds \emph{surely} at every
$\varphi\in\mathcal V$. \emph{\eqref{eq:normconst} does not move with $\mathsf L$, and the reason is which
bound on the contrast it reads}: the estimate behind it is taken at each contrast index separately, so it
consumes the \textbf{per-cell} bound $|\widehat{\widetilde W}_{i,l}|\le1$ of
\suppref{Remark}{hatTread}\textup{(2)} and not the Euclidean bound on the vector $\widehat{\widetilde W}_i$;
$\Pi_{\mathsf w}$ is a probability measure by B1(iv.d), so integrating out the third factor of $\Pi$
contributes the factor $1$ and $c_\Pi$ of B1(iv.c) is unchanged.
Run the iteration below on $\omega_n^{-1}\hat{\mathcal A}_\varphi$ and
$\omega_n^{-1}\hat T_{h_x}$, so that $\|\omega_n^{-1}\hat{\mathcal A}_\varphi\vert_{\mathcal S_n}\|\le1$
surely. By \suppref{Lemma}{scaleinv}(i)--(ii) this is a choice of relaxation parameter and changes neither the
iterates nor the stopping index; by (iii)--(iv) it leaves the rate \eqref{eq:phirate} unchanged; by (v)
it inflates the stopping index by $\omega_n^{2}$, and the iteration cap is sized accordingly,
\begin{equation}\label{eq:cap}
\bar m_n\;\ge\;\omega_n^{2}\,\hat\delta_n^{-2} .
\end{equation}
% ARCHAEOLOGY.  Through round 9 the normalisation was asserted to "cost nothing"
% and the cap was delta-hat_n^{-2}; eq:cap is what it actually costs, and it is a
% computation and not an assumption.

The stopping levels are computed before the data are touched. Fix $\mathsf c_{14},\mathsf c_{15}<\infty$
and $\kappa:=2\mu/(2\mu+1)$, and define the deterministic \emph{stagewise noise schedule}
\begin{equation}\label{eq:lfsched}
\hat\delta^{\mathrm{eff}}_{n,k}\;:=\;\mathsf c_{14}\Big(\hat\delta_n\ln n\;+\;h_u^{-2}\,\hat\varepsilon_{k-1}^{2}\Big),
\qquad
\hat\varepsilon_{k}\;:=\;\mathsf c_{15}\big(\hat\delta^{\mathrm{eff}}_{n,k}\big)^{\kappa},
\qquad k\ge1 .
\end{equation}
$\mathsf c_{14}$ and $\mathsf c_{15}$ must be
\emph{chosen}, because they are not identified by the data. What they must satisfy is this.
% ARCHAEOLOGY.  A round-9 draft said instead that every quantity in eq:lfsched was
% "known".
\Cref{prp:phirate} Step~2
(L3) needs $\mathsf c_{14}\ge\mathsf c_{16}$, with $\mathsf c_{16}$ the constant of \suppeq{effnoise};
Step~3 needs $\mathsf c_{15}\ge\mathsf c_{17}$, with $\mathsf c_{17}=2c_\mu\mathsf R_\ast^{1/(2\mu+1)}$
the constant of \suppeq{linlf}. Neither $\mathsf c_{16}$ nor $\mathsf c_{17}$ is known to the analyst:
the first is built from the kernel constants and the design, the second from (F-f)'s $c_\mu$, the
constant in the Landweber--Fridman bound under a discrepancy stop, which depends only on the source
index $\mu$ and on A7's $\tau_{\mathrm{DP}}$, and on nothing about the operator or the sample, 
and the \emph{spectral pre-image} radius $\mathsf R_\ast$ of (P-a.3). Which way the error bites, and why the
radius here is $\mathsf R_\ast$ and \textbf{not} $\mathsf R$, are in the supplement's discussion of the
estimator. For $k=1,\dots,\bar k$:
\begin{enumerate}
\item[(a)] form the linear equation in $v\in\mathcal S_n$,
$\ \hat{\mathcal A}_{\hat\varphi^{(k-1)}}v=-\hat T_{h_x}\big(\hat\varphi^{(k-1)},\theta\big)$;
\item[(b)] run Landweber--Fridman on it from $v_0=0$, at the relaxation \eqref{eq:normconst},
$v_{m+1}=v_m-\omega_n^{-2}\,\hat{\mathcal A}^\ast_{\hat\varphi^{(k-1)}}\big[\hat{\mathcal A}_{\hat\varphi^{(k-1)}}v_m
+\hat T_{h_x}(\hat\varphi^{(k-1)},\theta)\big]$, projected onto $\mathcal S_n$, which by
\suppref{Lemma}{scaleinv}(i) \emph{is} the unrelaxed iteration for the normalised pair;
\item[(c)] stop at
\[
\hat m^{(k)}\;:=\;\min\Big\{m\le\bar m_n:\ \big\|\hat{\mathcal A}_{\hat\varphi^{(k-1)}}v_m
+\hat T_{h_x}(\hat\varphi^{(k-1)},\theta)\big\|_{L^{2}(\Pi)}\;\le\;
\tau_{\mathrm{DP}}\,\hat\delta^{\mathrm{eff}}_{n,k}\Big\},
\]
the discrepancy principle \emph{at the stage-$k$ level} \eqref{eq:lfsched}, \emph{with the convention
that $\hat m^{(k)}:=\bar m_n$ when the set is empty};
\item[(d)] set $\hat\varphi^{(k)}:=\hat\varphi^{(k-1)}+v_{\hat m^{(k)}}$.
\end{enumerate}
Put $\hat\varphi:=\hat\varphi^{(\bar k)}$.

Write $\mathcal I_\ast$ for the fold the iteration is run on, $\mathcal I_{-k}$ under
\cref{def:crossfit}\textup{(b)}, $\mathcal I_m$ under \suppref{Definition}{hatTexp}, and read every ``measurable''
below at $\sigma(\mathcal I_\ast)$.
% ARCHAEOLOGY.  The letter was cal J through dispatch 21; the change is dispatch
% 22's.  cal J is the SINGULAR-SYSTEM INDEX SET of prp:compact's statement,
% cal J subset N, so one calligraphic letter carried an observation index here
% (i in cal J) and a singular-value index there (j in cal J), both in STATEMENTS,
% neither inside a proof.  The starred cal I is not a new letter: cal I is already
% this project's letter for a fold (cal I_a, cal I_m, cal I_k, cal I_{-k}), this
% object IS one of those two folds, and the star is the corpus's mark for the
% designated member of a family (cal T_star, cal G_star, Xi_star, sans-R_star,
% omega_star).  Glossary C81.

The output is a measurable random element, and that is \textbf{verified} rather than required
elsewhere. \Cref{def:crossfit}\textup{(c$'$.1)} needs each $\hat\varphi^{(-k)}$ to be a
$\sigma(\mathcal I_{-k})$-measurable random element of $\mathcal V$, and clause \textup{(b)} of that
definition names \emph{this} estimator as the object, so the property cannot be imposed there as a
construction requirement. \emph{One cannot require a property of an object one has already built}, which is
\suppref{Remark}{obligations}'s test read on an estimator rather than on an assumption. Induction on
$k\le\bar k$, from a $\sigma(\mathcal I_\ast)$-measurable $\hat\varphi^{(0)}$, gives the claim for
$\hat\varphi=\hat\varphi^{(\bar k)}$; \emph{the verification is three steps, each elementary and each
consuming a feature of the construction a reader could take for bookkeeping, followed by a fourth reading
that this paragraph's own consumers force}, and it is carried out in the supplement's discussion of the
estimator. \emph{Two of the features it consumes are \textup{(c)}'s}: the cap \eqref{eq:cap} is
deterministic, so the family of stopping events is \textbf{finite} and no limit is taken, and the emptiness
convention makes $\hat m^{(k)}$ \textbf{total}. \emph{Three stated results consume this paragraph}, each a
probability statement about $\hat\varphi$ and none of them well posed without it: \eqref{eq:phirate};
\suppref{Lemma}{pathcont}'s bound on $\max_{j\le\bar k}$ with probability tending to one; and (P-a.3)'s
event, which is defined through the realised operator at $\hat\varphi^{(k-1)}$.

Two features of \textup{(c)} are load-bearing, the level is stage-dependent, and it is
deterministic and supplied from outside, and both are argued in the supplement's discussion of the
estimator, where the status of Bakushinskii's veto is also fixed. The price actually paid is
separate from both. The schedule requires a \emph{rate}
$\hat\varepsilon_0$ for the preliminary and not mere consistency. \emph{That price is not new, only
newly visible}: the claim that $\bar k$ is finite and does not grow with $n$ already required
$\hat\varepsilon_0$ to be polynomially small. Hypothesis (P-a.2) of \cref{prp:phirate} carries it.

\emph{Every Landweber iteration in this construction is applied to a \textbf{linear} equation}; the
nonlinearity is carried by the outer index $k$ and by nothing else. That is the entire difference from
the route \suppref{Remark}{lfnonlinear} declines, and it is what makes a cone condition unnecessary.
\end{definition}

% ============================================================================
%  END OF PART H.
% ============================================================================

% ############################################################################
% ##  PART I -- memo Sec.10 STEP 7.  thm:rootn, at plan Sec.2.6's scope, with
% ##  cor:ratefloor (T2') as the floor.  Dispatch 10.
% ##
% ##  WHAT THIS PART SUPPLIES THAT NO EARLIER PART DID.
% ##   (1) THE MOMENT FUNCTION HAS A DEFINING DISPLAY.  psi is cited in Parts F
% ##       and G (lem:secondord, prp:orthog, prp:sharpquad) and was declared only
% ##       in notation_glossary.md.  def:moment fixes it in the .tex, together
% ##       with J and Sigma, which the memo's target statement names and which no
% ##       file displayed.
% ##   (2) THE ESTIMATOR.  def:crossfit and def:onestep.  def:lfstep defines the
% ##       varphi-step; nothing defined the theta-step.
% ##   (3) THE PRELIMINARY FOR theta.  lem:prelimtheta.  This is the object the
% ##       round-3 strategist-critic's m-23 named, and eq:prelimrate is NOT it --
% ##       see rem:rootnhyp(4).
% ##
% ##  ORDER OF WORK: def:moment, def:crossfit, def:onestep; lem:momzero;
% ##  lem:prelimtheta (L-14 with a rate); lem:biasexp (L-6 and L-7 instantiated
% ##  at the estimated nuisances); lem:epterm (L-8); lem:sandwich (L-11);
% ##  thm:rootn; cor:ratefloor (T2'); prp:medianpart (L-10); three remarks.
% ############################################################################

% ===========================================================================
%  I.1  THE MOMENT FUNCTION, THE JACOBIAN, THE VARIANCE
% ===========================================================================

\begin{definition}[The debiased moment function, and the two matrices]\label{def:moment}
For a nuisance triple $\zeta=(\varphi,L,\varrho)$ with $\varphi\in\mathcal V$,
$L:\IR\times\IR^{p}\to\IR^{p\times(\mathsf L-1)}$ Borel and
$\varrho:\IR^{p}\to\{q\in\IR^{\mathsf L-1}:q_l\ge0\ \text{for every }l,\ \mathbf 1'q\le1\}$ Borel, write
$\widetilde W_l(\varrho):=\un\{W=l\}-\varrho_l(\Delta Z)$ for $l=1,\dots,\mathsf L-1$ and
$\widetilde W(\varrho):=\big(\widetilde W_1(\varrho),\dots,\widetilde W_{\mathsf L-1}(\varrho)\big)'$, the
contrast vector \eqref{eq:contrast} formed at the trial propensity, and define the
$\IR^{p}$-valued \emph{debiased moment function}
\begin{equation}\label{eq:psidef}
\psi(\theta,\zeta)\;:=\;\Delta Z\,\Delta U(\theta,\varphi)\;+\;
L\big(\Delta U(\theta,\varphi),\,\Delta Z\big)\,\widetilde W(\varrho),
\end{equation}
with $\Delta U(\theta,\varphi)$ as in \eqref{eq:trialres}. Write $\widetilde W:=\widetilde W(\varrho)$ at the
true propensity, which is \eqref{eq:contrast} itself,
$\zeta_\star:=(\vdagger,L_\star,\varrho)$ as in \cref{def:weight}, and
\begin{equation}\label{eq:psistar}
\psi_\star\;:=\;\psi(\theta_\dagger,\zeta_\star)\;=\;\Delta Z\,\Delta U+L_\star(\Delta U,\Delta Z)\,\widetilde W.
\end{equation}
Define
\begin{equation}\label{eq:JSig}
J\;:=\;E\big[\Delta Z\,\Delta Z'\big]\;+\;E\big[\ell_\star(\Delta U,\Delta Z)\,\widetilde W\,\Delta Z'\big],
\qquad
\Sigma\;:=\;E\big[\psi_\star\psi_\star'\big] .
\end{equation}
\emph{Every product above is read at the type declaration of the instrument block}: $L$, $L_\star$ and
$\ell_\star$ carry one column per contrast, $\widetilde W(\varrho)$ is a column of $\mathsf L-1$ contrasts,
so $L\widetilde W(\varrho)$ is $\IR^{p}$-valued and the second summand of \eqref{eq:JSig} is
$(p\times1)(1\times p)$, the shape $J$ must have. \emph{At $\mathsf L=2$ every display above is the printed
one}: the simplex is the interval $[0,1]$, $\un\{W=1\}=W$ makes
$\widetilde W(\varrho)=W-\varrho(\Delta Z)$ a scalar, $L$, $L_\star$ and $\ell_\star$ have one column, and
$L\widetilde W=\widetilde WL$, $L_\star\widetilde W=\widetilde WL_\star$ and
$\ell_\star\widetilde W=\widetilde W\ell_\star$ for that reason.
% EXTRACTED FROM THE STATEMENT OF def:moment on 2026-08-03, under the house rule
% (.claude/references/theorem-environment-style.md) that a definition carries
% CONSTRUCTION ONLY.  Three justification blocks left; each is now a clause of
% rem:momentread, verbatim, in results_supp.tex (and results.tex).  NO \label
% MOVED: eq:psidef, eq:psistar and eq:JSig all stayed in the statement.
%   -- WHY THE SECOND SUMMAND CARRIES ell_star AND NOT L_star, why it is the zero
%      matrix under A1, and why the two-summand form is kept: rem:momentread(1).
%   -- THE "DIFFERENT MATRIX" CORRECTION RECORD, and the conditionality that
%      survives it:                                          rem:momentread(2).
%   -- THE TYPING NOTE, form against value, and its transfer to hat J:
%                                                            rem:momentread(3).
% ARCHAEOLOGY.  A sentence asserting the opposite stood here through dispatch 12:
% "J is the matrix of A3 and it is not E[Delta Z Delta Z']".  The claim entered as
% a round-2 CORRECTION -- round 1 had routed identification of theta_dagger through
% A3, and the repair (M-3) routed it through A1b(iv) instead, on the ground that A3
% is "nonsingularity of a different matrix" and is itself conditional on A2.  The
% conditionality half of that ground is right and survives; the "different matrix"
% half is false, and it was asserted and never computed.  Its negation was already
% proved in this same file at lem:nusigma(v).
\end{definition}

\noindent\emph{Notes on \cref{def:moment}.} Why the second summand of \eqref{eq:JSig} carries
$\ell_\star$ and not $L_\star$ and is nevertheless the zero matrix, why $J$ is not ``a different matrix''
from $E[\Delta Z\Delta Z']$, and how the typing of $E[\ell_\star\widetilde W\Delta Z']$ transfers to
$\hat J$ are \suppref{Remark}{momentread}.

% ===========================================================================
%  I.2  THE ESTIMATOR
% ===========================================================================

\begin{definition}[Cross-fitting, and the fixed anchor of the $\varphi$-step]\label{def:crossfit}
Fix $\mathsf F\ge2$, not depending on $n$, and a partition of $\{1,\dots,n\}$ into folds
$\mathcal I_1,\dots,\mathcal I_{\mathsf F}$ with $n_k:=|\mathcal I_k|$ and
$\lfloor n/\mathsf F\rfloor\le n_k\le\lceil n/\mathsf F\rceil$, chosen independently of the data. Write
$\mathcal I_{-k}:=\{1,\dots,n\}\setminus\mathcal I_k$, $O_i$ for the $i$-th observation vector of A0, and
\[
\IE_n[\varsigma]:=\frac1n\sum_{i=1}^{n}\varsigma(O_i),
\qquad
\IE_{n,k}[\varsigma]:=\frac1{n_k}\sum_{i\in\mathcal I_k}\varsigma(O_i) .
\]
At $\mathsf F=2$ this is the split of \cref{def:hatT}, with $(\mathcal I_k,\mathcal I_{-k})=(\mathcal I_m,\mathcal I_a)$.

\emph{Convention \textup{(P8)}, re-declared here.} \textup{(P8)}'s \emph{conditioning convention}, ``every expectation and every $O_P$ is
conditional on the auxiliary fold'', was stated \textbf{for the $\varphi$-step}, at $\sigma(\mathcal I_a)$, and is
therefore \textbf{re-declared} here rather than inherited: throughout the $\theta$-step it is read at
$\sigma(\mathcal I_{-k})$, with $\mathcal I_k$ in place of $\mathcal I_m$. Every conditional expectation and every conditional $O_P$ below is at that
$\sigma$-field, and A0 supplies the independence of $\mathcal I_k$ from $\mathcal I_{-k}$, which with
clause \textup{(c$'$)} below is what \textup{(P8)}'s \emph{freezing lemma} consumes.
\emph{Which clause of \textup{(P8)} each site uses is named at the site, and the answer at every site of
the $\theta$-step is \textbf{never the transfer rule}}; \suppref{Lemma}{pconv}(i)--(iii) are the three
clauses, and which inequality reaches which is in the supplement's discussion of the estimator.
% ARCHAEOLOGY.  Through dispatch 13 this paragraph said otherwise.

For each $k$ the nuisances are computed on $\mathcal I_{-k}$:
\begin{enumerate}
\item[(a)] \emph{Propensity.} $\hat\varrho^{(-k)}=(\hat\varrho^{(-k)}_1,\dots,\hat\varrho^{(-k)}_{\mathsf L-1})'$
satisfies A4(iii) with $n$ replaced by $n_{-k}$ and is
truncated to the $\varepsilon_\varrho$-interior simplex \eqref{eq:trunset} as in \cref{def:hatT}, so that
the fold contrast vector
\[
\widehat{\widetilde W}^{(-k)}\;:=\;\big(\un\{W=1\}-\hat\varrho^{(-k)}_1(\Delta Z),\ \dots,\
\un\{W=\mathsf L-1\}-\hat\varrho^{(-k)}_{\mathsf L-1}(\Delta Z)\big)' ,
\]
written $\widehat{\widetilde W}^{(-k)}_i$ when evaluated at observation $i$ and
$\widehat{\widetilde W}^{(-k)}_{i,l}$ for its $l$-th entry there, on the indexing of \cref{def:hatT},
satisfies $\varepsilon_\varrho\le\big|\widehat{\widetilde W}^{(-k)}_{i,l}\big|\le1$ surely \emph{at every}
$l\in\{1,\dots,\mathsf L-1\}$, this is \suppref{Remark}{hatTread}\textup{(2)} read at the fold, and it is
the \textbf{per-cell} bound and \textbf{not} a bound on $\|\widehat{\widetilde W}^{(-k)}_i\|$.
\emph{No statement \textbf{of this file} uses that Euclidean norm; several proofs in the supplement need
exactly it, and they \textbf{derive} it from the per-cell bound above rather than assume it}, 
\suppeq{estcontrnorm} is that derivation, and \suppref{Lemma}{epterm}, \suppref{Lemma}{sandwich} and
\suppref{Theorem}{rootn} are its consumers. So the sentence above scopes the absence to this file
and is not a statement that the object is unwanted. And
% MISDIRECTION CLOSED 2026-08-10 (Phase C).  This read "... and NOT a bound on
% ||What-tilde^{(-k)}_i||, WHICH NO STATEMENT IN THIS FILE USES".  Literally true --
% the scope is this file -- but eq:estcontrnorm derives precisely that norm bound
% from precisely this per-cell bound, and lem:sandwich Steps 2/5/9, lem:epterm
% Step 5 and thm:rootn Step 1 all consume it; proofs.tex's own tag at
% eq:estcontrnorm cites THIS clause by name and says the per-cell bound "is not a
% bound on ||What-tilde^{(-k)}_i||, which is exactly what an outer-product norm
% needs".  A reader who stopped here was told the object is unwanted.  The
% constant itself is deliberately NOT printed here: c_Wtilde is declared in
% results_supp.tex and proofs.tex and in no other file, and re-declaring it in a
% third would be an INV-7 collision.
$\mathsf r_k:=\hat\varrho^{(-k)}-\varrho$ satisfies $\|\mathsf r_k\|_\infty\le1-2\varepsilon_\varrho<1$
\textbf{by A4(i)}, the supremum taken over $z$ \emph{and} over the $\mathsf L-1$ components.
\emph{A4(i) is named because the \textbf{constant} rests on it and nothing else on this line does}:
\eqref{eq:trunset} puts $\hat\varrho^{(-k)}_l$ in $[\varepsilon_\varrho,1-\varepsilon_\varrho]$ by
construction, and it is A4(i), read at \emph{every} $l\in\{0,\dots,\mathsf L-1\}$, the base cell
included, that puts $\varrho_l$ in the same interval; without it $\varrho_l$ ranges over $[0,1]$ and
the constant is $1-\varepsilon_\varrho$, not $1-2\varepsilon_\varrho$. \emph{And this display is
\textbf{not} \suppeq{estcontrnorm}}: that one is \emph{Euclidean}, carries the constant declared there,
and is derived from the truncation alone with A4(i) expressly \emph{off} its list. They are a
sup-norm and a Euclidean statement about the same object, and only this one reads an overlap condition
, which is what makes them confusable and is why the difference is printed.
% A4(i) NAMED HERE 2026-08-10 (Phase 6a).  The bound was stated with NO ground: the
% constant 1 - 2 eps_varrho needs varrho_l in [eps, 1 - eps], which is A4(i), and
% this clause named neither.  THE ALTERNATIVE REPAIR -- CUTTING THE BOUND AS
% UNCONSUMED -- WAS CONSIDERED AND REJECTED, AND THE GROUND IS RECORDED BECAUSE THE
% CENSUS THAT SUPPORTED CUTTING WAS A STRING CENSUS.  "1 - 2 eps_varrho" does occur
% at exactly one site, its own; but proofs.tex's eq:estcontrnorm derivation cites
% this bound BY DESCRIPTION -- "def:crossfit(a)'s TWO per-cell bounds", the second
% of the two being this one -- which no search for the constant finds.  Cutting
% would have left a compiled proof citing a display that no longer exists, which is
% a worse defect than the unstated ground it removes.  The confusability with
% eq:estcontrnorm that motivated cutting is closed directly, by the sentence above.
% RECORDED FOR THE proofs.tex PASS, NOT DONE HERE: eq:estcontrnorm's sqrt(L-1) route
% for r does not actually need this bound -- varrho_l, hat-varrho_l in [0,1] already
% give |r_l| <= 1 -- so that derivation could be regrounded on the weaker fact and
% the phrase "two per-cell bounds" retyped.  That is a proofs.tex edit and is out of
% this pass's scope.  eq:trunset is declared in assumptions_main.tex, the first
% input, and this file already \eqref-s it twice, so the reference above adds no new
% dependency; the constant c_Wtilde is deliberately NOT named here, per the note
% ABOVE, which is why the sentence says "the constant declared there".
At $\mathsf L=2$,
\eqref{eq:trunset} is the interval $[\varepsilon_\varrho,1-\varepsilon_\varrho]$, $\un\{W=1\}=W$ makes
$\widehat{\widetilde W}^{(-k)}=W-\hat\varrho^{(-k)}(\Delta Z)$ a scalar, the cell subscript is vacuous, and
both bounds are the ones earlier rounds printed.
\item[(b)] \emph{Structural function.} $\hat\varphi^{(-k)}$ is the estimator of \cref{def:lfstep}, run on
$\mathcal I_{-k}$ \textbf{at a fixed anchor $\theta_\circ\in\IR^{p}$ chosen before the data are seen}, and
$h_k:=\hat\varphi^{(-k)}-\vdagger$. The anchor is what removes the circularity, and it is licensed and
not free. By \cref{cor:opsol} and \cref{lem:shiftinv} the population solution set of \eqref{eq:opeq} in
$\varphi$ is $\mathcal V_{CI}$ at \emph{every} $\theta$, so under A1b(iii) the $\varphi$-step targets
$\vdagger$ at $\theta_\circ$ as well as at $\theta_\dagger$; \emph{what is not $\theta$-free is the
conditioning of the inverse problem}, so \cref{prp:phirate}'s hypotheses (P-a)--(P-g), \eqref{eq:muhalf} and
A7 in the form \eqref{eq:a7amend} are read \emph{at $\theta_\circ$}, and $a$, $\mu$, $\mathsf R_\ast$ and
$\hat\delta_n$ are $\theta_\circ$-specific. This is \suppref{Remark}{R1}\textup{(iv)}, used and not extended.
\item[(c)] \emph{Debiasing weight.} $\hat L^{(-k)}$ is an admissible weight in the sense of
\cref{def:weight}(b), $\hat L^{(-k)}=L^{\,\hat b^{(-k)}}_{\hat\ell^{(-k)}}$, generated by an estimator
$\hat\ell^{(-k)}$ of $\ell_\star$ that is Borel, bounded, and continuous in $u$ at every $(u,z)$;
$D_k:=\hat L^{(-k)}-L_\star$. \emph{Those three properties of $\hat\ell^{(-k)}$ are the minimum
\cref{def:weight}(b) needs to be well posed}; hypothesis (R-a) of
\cref{thm:rootn} strengthens continuity to continuous differentiability, which \emph{one} clause of this
definition uses and only one, \textup{(c$'$.2)} needs it to make the inner supremum of the second
envelope, $\sup_u\|\partial_u\hat\ell^{(-k)}(u,z)\|$, attained on a countable dense set and hence
measurable. So the sup-norm convention \eqref{eq:supnormconv} is \textbf{not} unconditional.
% ARCHAEOLOGY.  Through dispatch 14 this sentence read "which no clause of this
% definition uses", from which a reader would conclude that it is unconditional.
\item[(c$'$)] Measurability of the five estimated nuisances, and the reading of their sup-norms.
\emph{Clause \textup{(c)} above states a per-realisation ``Borel'' and nothing more; without the present
clause the hypotheses of (R-a) and the displays \eqref{eq:foldrates} are not well-posed probability
statements, because they are statements about the distribution of a supremum and of an integral of a random
function.} \emph{And the right container is not the same for all five}, which is the substance of
\textup{(c$'$.1)}. \textup{(c$'$.1)} is a requirement on the estimators, with \textbf{one} exception. For $\hat\varphi^{(-k)}$ the property is \textbf{verified} rather than required,
at the measurability paragraph of \cref{def:lfstep}, which the supplement's discussion of the estimator
carries out.
\begin{enumerate}
\item[(c$'$.1)] Measurability, and it is two different statements, because three of the five are
only ever integrated and two are evaluated pointwise at arguments their container does not control.

The three integrated nuisances, as random elements of $L^{2}$ at a fixed representative. Each of
$\hat\varrho^{(-k)}$, $\hat b^{(-k)}$ and $\hat\varphi^{(-k)}$ is a
\textbf{$\sigma(\mathcal I_{-k})$-measurable random element} of, respectively, $L^{2}(P_{\Delta Z})$,
$L^{2}(P_{\Delta Z})$ and $\mathcal V$, \emph{a measurable map from the sample space into the space, and
not a family of functions indexed by the sample point}. A random element of an $L^{2}$ space over a
separable metric space admits a jointly measurable representative, measurable with respect to the
product of the sample $\sigma$-algebra with $\mathfrak B(\IR^{p})$ for the first two and with
$\mathfrak B(\mathcal X)$ for the third, \emph{fix one for each, once, as A2c does for
$\ell^\Pi_\star$}.
For the third of the three this is a \textbf{discharged verification} and not a requirement, and the
difference is not cosmetic. Clause \textup{(b)} pins $\hat\varphi^{(-k)}$ to \cref{def:lfstep}, whose
stopping index is \emph{random}, so the ground that settles the other two, A4(iii)'s ``by inspection of the
formula'', is unavailable, and what supplies it instead is that definition's measurability paragraph.
\emph{That the choice of representative is immaterial for these three is the test a container has to pass};
it is passed on two different grounds, and both are in the supplement's discussion of the estimator.

For $\hat\varrho^{(-k)}$ the choice is moreover \textbf{pinned}, and that is what clause
\textup{(a)}'s \emph{sure} bound rests on. A4(iii), read at the fold by
clause \textup{(a)}, requires product measurability of $\hat\varrho$'s \emph{own pointwise function}, 
the \eqref{eq:prodmeas} shape, so for that object no equivalence class is formed and
no representative is chosen at all, and
$\varepsilon_\varrho\le\big|\widehat{\widetilde W}^{(-k)}_{i,l}\big|\le1$ holds \emph{surely} at every $i$
and every $l$, as clause \textup{(a)} says.
% ARCHAEOLOGY.  Through dispatch 15, clause (a) asserted a SURE property of an
% object this clause had placed in an L^2 container, with nothing tying the chosen
% representative to the truncation -- this file's own test, applied to the wrong
% object.  Every consumer of clause (a) would have survived a weakening to "almost
% surely", so it was a wording defect and not a broken bound; it was repaired at
% A4(iii) in dispatch 16 because the assumption had to be repaired anyway.

The two weight objects, for which the $L^{2}$ container is not merely unnecessary but wrong. What is
required of $\hat\ell^{(-k)}$, in addition to clause \textup{(c)}, is that its \emph{own} pointwise
function, the one \textup{(c)} already defines, at every realisation and every $(u,z)$, be product
measurable:
\begin{equation}\label{eq:prodmeas}
(\omega,u,z)\;\longmapsto\;\hat\ell^{(-k)}(\omega)(u,z)
\quad\text{is}\quad
\sigma(\mathcal I_{-k})\otimes\mathfrak B\big(\IR\times\IR^{p}\big)\text{-measurable} ,
\end{equation}
and $\hat L^{(-k)}=L^{\,\hat b^{(-k)}}_{\hat\ell^{(-k)}}$ then \emph{inherits} the same property from
\eqref{eq:weightdef}, together with the representative fixed above for $\hat b^{(-k)}$, \emph{so
\eqref{eq:prodmeas} is imposed on one object and not on two}. \emph{Under (R-a) the derivative inherits it too, and \textup{(c$'$.2)} needs that for its
second envelope}. \emph{The $L^{2}$ reading fails for these two and not for the three above} because they are
evaluated at a \emph{shifted} first argument whose law \textup{(c$'$.2)} needs \emph{not} to be dominated by
$P_{\Delta U,\Delta Z}$, and on an undominated set two representatives of one class disagree; both
inheritances are verified in the supplement's discussion of the estimator.

\emph{What is retained from the $L^{2}$ reading}, because several results use it, is membership:
$\hat\ell^{(-k)}\in L^{2}(P_{\Delta U,\Delta Z})$ almost surely, \suppref{Lemma}{balance} Step~0's third
point, and what makes the third display of \eqref{eq:foldrates} finite, and, for $\hat L^{(-k)}$, the
same membership, which is \emph{derived and not required}. \emph{The third display of \eqref{eq:foldrates} is a
statement with a truth value}: by \eqref{eq:prodmeas} and Tonelli's theorem,
$\omega\mapsto\|\hat\ell^{(-k)}-\ell_\star\|_{L^{2}(P)}$ is $\sigma(\mathcal I_{-k})$-measurable, and
likewise every conditional expectation and every $O_P$ statement below about these objects is well posed.
Equation \eqref{eq:prodmeas} is also what \textup{(P8)}'s freezing lemma consumes for the
$\hat\ell^{(-k)}$ \emph{ingredient} of its composite.
\item[(c$'$.2)] The sup-norms are essential suprema in $z$ and genuine suprema in $u$. For the two
envelopes (R-a) prices,
\begin{equation}\label{eq:supnormconv}
\big\|\hat\ell^{(-k)}\big\|_\infty
\;:=\;\operatorname{ess\,sup}_{z}\ \Big\{\ \sup_{u\in\IR}\ \big\|\hat\ell^{(-k)}(u,z)\big\|\ \Big\} ,
\qquad\text{the essential supremum taken under }P_{\Delta Z},
\end{equation}
and likewise for $\|\partial_u\hat\ell^{(-k)}\|_\infty$. Both halves of that reading are needed and
neither can be interchanged with the other. The inner supremum is over an uncountable index but is
attained on a countable dense set, because \textup{(c)} makes $u\mapsto\hat\ell^{(-k)}(u,z)$ continuous at
every $z$ and (R-a) makes $u\mapsto\partial_u\hat\ell^{(-k)}(u,z)$ continuous as well, so
$(\omega,z)\mapsto\sup_u\|\cdot\|$ is jointly measurable by \textup{(c$'$.1)}. \emph{The \textbf{outer}
essential supremum is measurable as well, and that does \textbf{not} follow from joint measurability without
an argument}, it is not a countable supremum, and the argument is in the supplement's discussion of the
estimator.
Hence $\omega\mapsto\|\hat\ell^{(-k)}\|_\infty$, and likewise
$\omega\mapsto\|\partial_u\hat\ell^{(-k)}\|_\infty$, is $\sigma(\mathcal I_{-k})$-measurable as a map into
$[0,+\infty]$. \emph{This is load-bearing and not hygiene}: (R-a) prices both envelopes as $O_P(1)$, which
is a statement about the distribution of these two maps.

\emph{And the essential supremum must be taken in $z$ and not
in $(u,z)$}, because every consumer of \eqref{eq:supnormconv} evaluates the envelope at a
\emph{shifted} first argument, whose law need not be dominated by $P_{\Delta U,\Delta Z}$, while its
\emph{second} argument is in every case a draw of $\Delta Z$. \emph{The shift takes two shapes, and the
list that follows is \textbf{illustrative} and not a complete enumeration of the consumers}:
$(\Delta U-\lambda\Delta h_k,\Delta Z)$ in \suppref{Lemma}{biasexp} Steps~1 and~4--5, and
$(\Delta U_i(\theta,\hat\varphi^{(-k)}),\Delta Z_i)$ in \suppref{Lemma}{sandwich} Steps~2 and~9.
\suppref{Lemma}{epterm} Steps~3 and~5 read the first shape; \suppref{Lemma}{sandwich} Steps~3 and~5 and
\suppref{Theorem}{rootn} Step~1 read the second.
% AMBIGUITY RESOLVED 2026-08-10 (Phase C).  The list read as an em-dash aside
% inside "every consumer ... evaluates the envelope at a shifted first argument",
% so its grammar was ambiguous between exhaustive and illustrative, and read as
% exhaustive it was short by five: lem:epterm Steps 3 and 5, lem:sandwich Steps 3
% and 5, and thm:rootn Step 1.  ALL FIVE VERIFIED AT THE ARTIFACT before this
% sentence was written, and each was checked to read a shift of one of the two
% shapes already named -- epterm at (Delta U - Delta h, Delta Z), sandwich 3/5 and
% rootn 1 at Delta U(theta, phi-hat^{(-k)}).  The list is now declared
% illustrative and the five are named, rather than a fresh exhaustiveness claim
% being made: the universal quantifier that carries the argument is "every
% consumer evaluates at a SHIFTED first argument", and that is what is asserted.
\end{enumerate}
\item[(d)] The fold reading of the maintained nuisance-rate clauses, declared once. A4(iii) is read
at each fold in (a); the \textbf{four} clauses of A5 that price an \emph{estimated} nuisance are read the
same way, and every later citation of them is a citation of the corresponding display here:
\begin{equation}\label{eq:foldrates}
\max_{k\le\mathsf F}\big\|h_k\big\|_{L^{2}(\pi)}=O_P(\delta_n),
\qquad
\max_{k\le\mathsf F}\big\|h_k\big\|_{H^{s}}=O_P(1),
\qquad
\max_{k\le\mathsf F}\big\|\hat\ell^{(-k)}-\ell_\star\big\|_{L^{2}(P)}=O_P(\epsilon_n),
\end{equation}
being A5(ii)'s first half, A5(vi), and A5(iv)'s first half respectively, \emph{the third in force precisely
when (R-e.1) of \cref{thm:rootn} is the branch taken, since under (R-e.2) no rate for $\hat\ell^{(-k)}$ is
assumed at all}. \emph{The second display is finite at every realisation only under
\eqref{eq:sieveinhs}, and it is cited here for the same reason A5(vi) carries it}: $h_k$ is
$\hat\varphi^{(-k)}-\vdagger$ with $\hat\varphi^{(-k)}$ an iterate of \cref{def:lfstep} and hence, 
\textbf{once $\hat\varphi^{(0)}$ is in $\mathcal S_n$}, an
element of $\mathcal S_n$, so A8's basis clause is what places it in $H^{s}$ and H1 is what places
$\vdagger$ there.
\emph{The premise in bold is stated here because this is a \textbf{definition} and inherits no hypothesis
list, and because the display does not hold without it}: \cref{def:lfstep}\textup{(d)} sets
$\hat\varphi^{(k)}=\hat\varphi^{(k-1)}+v_{\hat m^{(k)}}$ with $v_m\in\mathcal S_n$, so the path lies in
$\mathcal S_n$ if and only if its \emph{starting point} does, which is why A5(vi$'$), (P-a.1) and
\cref{def:lfstep}'s own opening all say ``once $\hat\varphi^{(0)}$ does''.
\emph{Every result that reads \eqref{eq:foldrates} supplies the premise},
(P-a.1) being in group (T-5) of \cref{thm:rootn}.
% ARCHAEOLOGY.  Through dispatch 20 this clause stated the display without the
% bold premise.
The indexing is not cosmetic
and it is not free either. A5(ii), (iv) and (vi) are stated at a \textbf{single} estimated nuisance, so a
result that cites them bare while averaging over $\mathsf F$ fold complements is short by exactly the
maximum. For the first two the cost is nil: \cref{thm:rootn} Step~0 \emph{derives} them from (T-5) with
$\mathsf F$ fixed and $n_{-k}\asymp n$. \emph{For the third it is a strengthening of A5(iv) in form},
of the same order for fixed $\mathsf F$ but a strengthening nonetheless, and it is therefore carried
explicitly by hypothesis (R-e.1) of \cref{thm:rootn} rather than read into A5(iv). (R-e.2) and (R-a)
already write $\max_{k}$.
\emph{The fourth clause is A5(vi$'$).} It prices the \emph{stagewise} Sobolev stability of the
$\varphi$-step,
% ARCHAEOLOGY.  (R-e.1) was left bare by a round-10 draft; this enumeration read
% "the three clauses" through dispatch 13.
$P(\max_{j\le\bar k}\|\hat\varphi^{(j)}-\vdagger\|_{H^{s}}\le\mathsf R)\to1$, where $\bar k$ is a
\emph{stage} index and not a fold index, the two being one glyph apart, and it is a clause about an
\textbf{estimated} nuisance, so it is read at each fold complement exactly as the other three are: for each
$k\le\mathsf F$, A5(vi$'$)'s display is required of the $\varphi$-step run on $\mathcal I_{-k}$.
That this costs nothing, and that the radius $\mathsf R$ must nonetheless serve all $\mathsf F$ fold
complements at one value, are in the supplement's discussion of the estimator.
\end{enumerate}
Put $\hat\zeta^{(-k)}:=(\hat\varphi^{(-k)},\hat L^{(-k)},\hat\varrho^{(-k)})$ and, for $\theta\in\IR^{p}$,
\begin{equation}\label{eq:barpsi}
\bar\psi_n(\theta)\;:=\;\frac1n\sum_{k=1}^{\mathsf F}\sum_{i\in\mathcal I_k}\psi\big(\theta,\hat\zeta^{(-k)}\big)(O_i),
\qquad
\hat\psi^{(-k)}:=\psi\big(\theta_\dagger,\hat\zeta^{(-k)}\big).
\end{equation}
\emph{Since $\mathsf F$ is fixed, $n_k/n\to\mathsf F^{-1}$ and $n_k\asymp n$}; the constant $\mathsf F$ is
carried explicitly in every bound rather than absorbed, because \suppref{Proposition}{medianpart} quantifies over
partitions and a suppressed $\mathsf F$ would hide the one place a partition enters an order.
\end{definition}

\begin{definition}[The preliminary and the one-step estimator]\label{def:onestep}
With the notation of \cref{def:crossfit} define
\begin{equation}\label{eq:prelimtheta}
\hat\theta^{(0)}\;:=\;\Big(\IE_n\big[\Delta Z\Delta Z'\big]\Big)^{-1}
\frac1n\sum_{k=1}^{\mathsf F}\sum_{i\in\mathcal I_k}\Delta Z_i\big(\Delta Y_i-\Delta\hat\varphi^{(-k)}_i\big),
\end{equation}
on the event that $\IE_n[\Delta Z\Delta Z']$ is nonsingular, and $\hat\theta^{(0)}:=0$ otherwise. Define
the \emph{cross-fitted Jacobian}, for $\theta\in\IR^{p}$,
\begin{equation}\label{eq:Jhat}
\hat J(\theta)\;:=\;\frac1n\sum_{k=1}^{\mathsf F}\sum_{i\in\mathcal I_k}
\Big[\Delta Z_i\Delta Z_i'+
\hat\ell^{(-k)}\big(\Delta U_i(\theta,\hat\varphi^{(-k)}),\Delta Z_i\big)\,
\widehat{\widetilde W}^{(-k)}_i\,\Delta Z_i'\Big],
\end{equation}
so that $\partial_\theta\bar\psi_n(\theta)=-\hat J(\theta)$ identically, and put
$\hat J:=\hat J(\hat\theta^{(0)})$. \emph{The order of the three factors is what the identity forces}: by
\eqref{eq:psidef} the $\theta$-derivative of $\hat L^{(-k)}(\Delta U_i(\theta,\cdot),\Delta Z_i)\,
\widehat{\widetilde W}^{(-k)}_i$ is $-\hat\ell^{(-k)}(\cdot)\,\widehat{\widetilde W}^{(-k)}_i\,\Delta Z_i'$,
which is $(p\times(\mathsf L-1))((\mathsf L-1)\times1)(1\times p)$. At $\mathsf L=2$ the three factors
commute in the first two positions and \eqref{eq:Jhat} is the display earlier rounds printed.
% WHY THE IDENTITY ABOVE USES ONE DERIVATIVE AND NO MORE, and why the second
% derivative is hypothesised at thm:rootn (R-a) and not here: rem:onestepread(1).
% The IDENTITY ITSELF STAYS -- only its proof moved.
The \emph{one-step estimator} is
\begin{equation}\label{eq:onestep}
\hat\theta\;:=\;\hat\theta^{(0)}\;+\;\hat J^{-1}\,\bar\psi_n\big(\hat\theta^{(0)}\big)
\end{equation}
on the event $\{\hat J\ \text{nonsingular}\}$, and $\hat\theta:=\hat\theta^{(0)}$ otherwise. The variance
estimator is
\begin{equation}\label{eq:Sighat}
\hat\Sigma\;:=\;\frac1n\sum_{k=1}^{\mathsf F}\sum_{i\in\mathcal I_k}
\psi\big(\hat\theta,\hat\zeta^{(-k)}\big)(O_i)\;\psi\big(\hat\theta,\hat\zeta^{(-k)}\big)(O_i)' ,
\qquad
\hat V\;:=\;\hat J^{-1}\hat\Sigma\,\hat J^{-1\prime} ,
\end{equation}
the last \textbf{on the event $\{\hat J\ \text{nonsingular}\}$, with $\hat V:=0$ otherwise}.
% EXTRACTED FROM THE STATEMENT OF def:onestep on 2026-08-03, under the house rule
% (.claude/references/theorem-environment-style.md) that a definition carries
% CONSTRUCTION ONLY.  Each block is now a clause of rem:onestepread, immediately
% below, verbatim.  NO \label MOVED: eq:prelimtheta, eq:Jhat, eq:onestep and
% eq:Sighat all stayed in the statement, and so did the one-derivative identity.
%   -- WHY THE FALLBACK IS NOT OPTIONAL, why the zero matrix is the value chosen,
%      and why it must not be reported as a variance: rem:onestepread(2).
%   -- ONE STEP AND NOT A ROOT, the sample-against-population ground, and the
%      literature anchors:                            rem:onestepread(3).
% THE REMARK IS IN THIS FILE AND NOT IN results_supp.tex, ON PURPOSE.  Clause (3)
% carries chernozhukov2018dml, newey1994h and vandervaart1998, and those are the
% ONLY citations of the last two anywhere in the corpus.  results_s3_estimation is
% read by BOTH main.qmd and supplement.qmd; results_supp.tex by the supplement
% only.  Routing the clause across documents would drop two anchors out of the
% article's reference list, so the block stays in the document it started in.
% ARCHAEOLOGY.  Through dispatch 12 clause (3)'s paragraph gave the population half
% of the ground as "A3 is a LOCAL rank condition ... and no global rank condition
% is available".
\end{definition}

\begin{remark}[How \cref{def:onestep}'s construction is read]\label{rem:onestepread}
Three clauses, each of which stood inside the definition. None of it is construction, and the
estimator \cref{def:onestep} defines is unchanged by its removal.
\begin{enumerate}
\item[(1)] \emph{That identity uses \textbf{one} derivative and no more}:
$\partial_u\hat L^{(-k)}=\hat\ell^{(-k)}$, which \cref{def:weight}(b) supplies at every point under the three
properties \cref{def:crossfit}(c) requires of $\hat\ell^{(-k)}$, together with $\theta\mapsto\Delta U_i(\theta,\cdot)$
being affine by \eqref{eq:trialres} and the sum being finite. No second derivative is claimed here.
$\partial_u^{2}\hat L^{(-k)}=\partial_u\hat\ell^{(-k)}$ is used at \cref{thm:rootn} Step~1, for Taylor's
theorem in $\theta$, and it is hypothesised \emph{there}, at (R-a), \emph{a definition should not carry a
differentiability claim that a theorem assumes}. \emph{That step is not the only reader of the
envelope $\|\partial_u\hat\ell^{(-k)}\|_\infty$, and no claim of exhaustiveness is made here}:
\suppref{Lemma}{biasexp} Step~4 and \suppref{Lemma}{sandwich} Step~5 read it too, on the right of an
inequality, as (R-a)'s own reading records. The draft this sentence replaces said ``only'' and named
one of three, while that reading named two of the same three; the set is one and the two claims
covered disjoint halves of it.
\item[(2)] \emph{The fallback is not optional.} $\hat\theta^{(0)}$ is given a value off $\{\IE_n[\Delta Z\Delta Z']\
\text{nonsingular}\}$ and $\hat\theta$ is given one off $\{\hat J\ \text{nonsingular}\}$;
$\hat V$ inverts the \emph{same} $\hat J$ and needs one too. That matters because
\cref{thm:rootn}'s \textbf{second} conclusion, $\hat V\xrightarrow{p}J^{-1}\Sigma J^{-1\prime}$, and the
studentised statement that follows it are statements about $\hat V$ at every $n$, and a convergence in
probability is not well posed for a quantity that is undefined on an event of positive probability.
\emph{The choice of value is immaterial and the zero matrix is chosen because it is the one value that
needs no further object}: by \suppref{Lemma}{sandwich}\textup{(i)}, $P(\hat J\ \text{nonsingular})\to1$, so the
event on which the fallback is active has probability tending to zero and no conclusion of
\suppref{Lemma}{sandwich}\textup{(ii)} or of \cref{thm:rootn} changes. \emph{In a finite sample the fallback is
not innocuous and should not be reported as a variance estimate}: $\hat V=0$ makes the studentised
statistic undefined, which is the honest report on that event and is the reason the value is stated rather
than left to the reader.
\item[(3)] One step and not a root, and the ground is the \textbf{sample} criterion and not the
population one. The realised map
$\theta\mapsto\bar\psi_n(\theta)$ is \emph{not} affine: by \eqref{eq:psidef} $\theta$ enters
$\hat L^{(-k)}$ through its first argument, and $\widehat{\widetilde W}^{(-k)}$ is not conditionally
centred, so $\bar\psi_n(\theta)=0$ need not have a unique root at any $n$. \emph{The ground is
\textbf{not} that no global rank condition is available}: by \suppref{Lemma}{momzero}\textup{(iv)} and
\suppeq{psiaffine} the
population moment $\theta\mapsto E[\psi(\theta,\zeta_\star)]$ is affine with gradient $-E[\Delta Z\Delta Z']$,
so a global rank condition \emph{is} available, it is A1b(iv), and $\theta_\dagger$ is its unique root on
all of $\IR^{p}$, which says nothing about the sample criterion at estimated nuisances.
\emph{What the block supplies is a preliminary with a
rate} (\suppref{Lemma}{prelimtheta}), and a one-step from a preliminary with a rate is exactly the construction
that converts a rank condition on the population Jacobian into a limit distribution without solving
anything; the cross-fitted form is that of
\citet{chernozhukov2018dml}, and the two classical anchors for the one-step itself are
\citet[\S3.4]{newey1994h} and \citet[Ch.~5]{vandervaart1998}. \emph{The Jacobian is
evaluated at $\hat\theta^{(0)}$ and not at $\hat\theta$}, which is what makes \eqref{eq:onestep} a closed-form
estimator and what makes the algebra of \cref{thm:rootn} Step~1 exact rather than asymptotic.
\end{enumerate}
\end{remark}
